% 2>/dev/null; MSYS_NO_PATHCONV=1 exec docker run --rm -v .:/work -w /work texlive/texlive:latest latexmk -pdf -interaction=nonstopmode -outdir="${0%.tex}" "$0"
% paper/ に置いた資産(図 figNN.png・refs.bib 等)は全てパス無しで参照できる(-outdir が入力探索先も兼ねるため)
\documentclass{article}

\usepackage[utf8]{inputenc}
\usepackage[T1]{fontenc}
\usepackage{amsmath}
\usepackage{amssymb}
\usepackage{graphicx}
\usepackage{booktabs}
\usepackage{hyperref}

% TODO: replace handle with the author's real name and affiliation before submission.
\title{An open-source pipeline for coupled MHD pressure-drop and tritium-breeding
analysis of liquid-metal flow channels on stellarator magnetic surfaces}
\author{lzpel}
\date{\today}

\newcommand{\Ha}{\mathrm{Ha}}

\begin{document}
\maketitle

\begin{center}
\fbox{\textbf{[DRAFT --- all numerical results and figures in this version are mock placeholders]}}
\end{center}

\begin{abstract}
Liquid-metal breeding blankets are a leading candidate for tritium self-sufficiency
in stellarator fusion power plants, but no public tool exists to evaluate whether a
given flow-channel layout on a three-dimensionally twisted magnetic geometry is
feasible: design methods established for tokamaks do not transfer, and stellarator
studies remain internal to individual laboratories. We present \texttt{mhd-tbr-stell},
an open-source pipeline that starts from a VMEC equilibrium (\texttt{wout.nc}),
generates liquid-metal flow channels conforming to offset magnetic surfaces, and
returns the three plant-level feasibility figures of a breeding blanket: the tritium
breeding ratio (TBR) computed with OpenMC on CAD geometry, the total
magnetohydrodynamic (MHD) pressure drop and the corresponding pumping-power fraction
computed with a correlation-based one-dimensional duct network, and the peak
structure temperature. The correlation network is anchored by fully resolved MHD
simulations (epotFoam/OpenFOAM) validated against the analytic Hartmann, Shercliff,
and Hunt solutions. As a demonstration we map the pressure-drop distribution of a
dual-coolant lead--lithium (DCLL) module sector on a quasi-isodynamic
(W7-X--like) configuration and trace the Pareto front between TBR and pumping-power
fraction. \textbf{The present manuscript is a structural draft: every quantitative
value is a placeholder to be replaced by production runs.}
\end{abstract}

\section{Introduction}

Tritium self-sufficiency requires the breeding blanket to achieve a tritium breeding
ratio (TBR) above approximately 1.1 while keeping the power needed to circulate the
breeder within a small fraction of the fusion power. For liquid-metal breeders such
as PbLi, the dominant penalty is the magnetohydrodynamic (MHD) pressure drop induced
by the confining field~\cite{smolentsev2010}. In tokamaks, both the neutronics and
the MHD sides of this trade-off have been studied for decades and the design
methodology is mature. Stellarators break this methodology: the blanket must follow
magnetic surfaces that twist in three dimensions, the field at the channel varies in
magnitude \emph{and} direction along the flow path, and the space between plasma and
coils is strongly non-uniform. Recent stellarator power-plant concepts --- among them
the quasi-isodynamic Stellaris design~\cite{lion2025stellaris} and the Infinity Two
pilot-plant study~\cite{infinitytwo2025} --- confirm that blanket feasibility is a
first-order design driver, yet the flow-path-level analysis behind such concepts
remains internal to each organisation.

On the geometry side, the open-source ParaStell package~\cite{moreno2024parastell}
generates parametric in-vessel components from a VMEC
equilibrium~\cite{hirshman1983} and couples them to the Monte Carlo transport code
OpenMC~\cite{romano2015} for neutronics. What is missing downstream is the liquid
metal itself: no public tool couples the MHD pressure drop of the flow channels to
the TBR of the same geometry. This paper describes \texttt{mhd-tbr-stell}, an
open-source pipeline built deliberately by \emph{plumbing existing, validated codes
together} rather than by writing a new solver, and applying them to a computation
that has not previously been published: the coupled MHD--neutronics feasibility map
of liquid-metal channels on stellarator magnetic surfaces.

Our contributions are:
(i) a channel generator that produces watertight CAD solids and body-fitted
structured meshes directly from the Fourier representation of a VMEC equilibrium;
(ii) a correlation-based one-dimensional MHD network that returns the total pressure
drop and pumping-power fraction of a full module sector in minutes, anchored by
resolved epotFoam simulations validated against analytic
solutions~\cite{hartmann1937,shercliff1953,hunt1965};
(iii) the first published TBR--pressure-drop Pareto front for a stellarator DCLL
module sector.

\section{Method}

\subsection{Channel geometry from a VMEC equilibrium}

The last closed flux surface is given by the VMEC Fourier series
\begin{equation}
R(\theta,\phi) = \sum_{m,n} R_{mn}\cos(m\theta - nN_{\mathrm{fp}}\phi),\qquad
Z(\theta,\phi) = \sum_{m,n} Z_{mn}\sin(m\theta - nN_{\mathrm{fp}}\phi),
\end{equation}
with $N_{\mathrm{fp}}$ field periods. Closed-form partial derivatives in $\theta$
and $\phi$ give surface normals and tangents analytically; channels are constructed
between two normal-offset surfaces and partitioned into poloidal--toroidal ducts and
feeding manifolds. Two discretisations are exported from the same parametric
description: (a) watertight B-rep solids with per-solid material labels, written as
STEP for CAD-based neutronics, and (b) body-fitted structured hexahedral meshes with
wall-normal grading sufficient to resolve Hartmann layers, written as OpenFOAM
\texttt{polyMesh}. The B-rep stage uses the cadrum kernel (statically linked
OpenCASCADE, Rust)~\cite{cadrum}.

\subsection{Fast MHD pressure-drop network}

For each duct segment the local field $\mathbf{B}$ from the equilibrium is projected
onto the duct cross-section. For fully developed flow in a rectangular duct at high
Hartmann number
$\Ha = B\,L\sqrt{\sigma/(\rho\nu)}$,
the pressure gradient is estimated with wall-conductance correlations of the
Miyazaki form~\cite{miyazaki1986}
\begin{equation}
\frac{\mathrm{d}p}{\mathrm{d}x} \simeq k_p\,\sigma\,u\,B^2,\qquad
k_p = \frac{c}{1+a/(3b)+c},\qquad
c = \frac{\sigma_w t_w}{\sigma L},
\end{equation}
with duct half-widths $a,b$, wall thickness $t_w$, and conductivities
$\sigma,\sigma_w$. Bends, expansions, and manifolds receive three-dimensional loss
coefficients from the literature~\cite{smolentsev2010}. Solving the resulting
one-dimensional network for a prescribed mass-flow distribution yields the total
pressure drop $\Delta p$ and the pumping-power fraction
$f_{\mathrm{pump}} = \dot{V}\Delta p / (\eta_{\mathrm{pump}} P_{\mathrm{fus}})$
in seconds, which is what permits design-space scans.

\subsection{Anchoring CFD and neutronics}

Representative straight and bent segments are recomputed with epotFoam, an
OpenFOAM solver for inductionless MHD based on the electric-potential
formulation~\cite{tassone2019epotfoam}, on the structured meshes of
Sec.~2.1. The solver is first validated against the analytic Hartmann, Shercliff,
and Hunt solutions (Sec.~3). The TBR is computed with OpenMC~\cite{romano2015} on
the DAGMC faceted representation of the same STEP geometry, following the
ParaStell workflow~\cite{moreno2024parastell}.

\section{Validation (mock)}

Table~\ref{tab:validation} summarises the reproduction of the canonical analytic
MHD duct flows. \textbf{All error values are placeholders.}

\begin{table}[htbp]
\centering
\caption{Validation of the anchoring CFD against analytic solutions
(\textbf{mock values}).}
\label{tab:validation}
\begin{tabular}{llrr}
\toprule
Case & Configuration & $\Ha$ & Error in $\Delta p$ \\
\midrule
Hartmann~\cite{hartmann1937} & parallel plates, uniform $B$ & $10^3$ & 0.8\,\% \\
Shercliff~\cite{shercliff1953} & rectangular duct, insulating walls & $10^3$ & 1.4\,\% \\
Hunt~\cite{hunt1965} & rectangular duct, conducting Hartmann walls & $10^3$ & 2.1\,\% \\
\bottomrule
\end{tabular}
\end{table}

Figure~\ref{fig:hartmann} compares the computed Hartmann velocity profile with the
analytic solution $u(y)/u_0 = \left(\cosh\Ha - \cosh(\Ha\,y/L)\right) /
\left(\cosh\Ha - \sinh(\Ha)/\Ha\right)$.

\begin{figure}[htbp]
\centering
\includegraphics[width=0.62\linewidth]{fig02.png}
\caption{Hartmann-flow velocity profile: analytic solution (line) versus anchoring
CFD (markers). \textbf{Mock figure: CFD markers are synthetic.}}
\label{fig:hartmann}
\end{figure}

\section{Results (mock)}

\subsection{Pressure-drop map of a DCLL module sector}

Figure~\ref{fig:dpmap} shows the distribution of the MHD pressure gradient over the
unrolled $(\theta,\phi)$ blanket surface of one field period of a W7-X--like
configuration, for a poloidal--radial DCLL channel layout at
$\Ha \approx 10^3$. High-loss regions correlate with the high-field, strongly
shaped bean cross-sections. \textbf{The map is a synthetic placeholder.} The total
sector pressure drop is $\Delta p = 0.87\,$MPa and the pumping-power fraction
$f_{\mathrm{pump}} = 1.2\,\%$ (\textbf{mock values}).

\begin{figure}[htbp]
\centering
\includegraphics[width=0.8\linewidth]{fig01.png}
\caption{MHD pressure-gradient distribution on the unrolled $(\theta,\phi)$ blanket
surface, one field period. \textbf{Mock figure.}}
\label{fig:dpmap}
\end{figure}

\subsection{TBR--pumping-power Pareto front}

Varying channel count, duct aspect ratio, breeder-zone thickness, and flow-channel
insert conductance produces the design population of Fig.~\ref{fig:pareto}. The
front exhibits the expected trade-off: thicker breeder zones and slower, wider ducts
raise the TBR while lowering $f_{\mathrm{pump}}$, until neutron shielding of the
coils bounds the available radial build. \textbf{All points are synthetic.}

\begin{figure}[htbp]
\centering
\includegraphics[width=0.62\linewidth]{fig03.png}
\caption{Pareto front between TBR and pumping-power fraction for the sampled
channel layouts. \textbf{Mock figure.}}
\label{fig:pareto}
\end{figure}

\section{Discussion}

Three limitations bound the present approach. First, the correlation network assumes
fully developed flow between discrete loss elements; strongly three-dimensional
effects (field-gradient regions, electrically coupled adjacent ducts) are captured
only through the anchoring CFD, and the fidelity of the fast branch is therefore
bounded by the segment library. Second, the anchoring simulations are performed at
$\Ha \sim 10^3$, one order below reactor conditions, because Hartmann-layer
resolution at $\Ha \sim 10^4$ requires meshes of order $10^8$ cells; extrapolation
relies on the $\Ha$-scaling built into the correlations. Third, buoyancy-driven
mixed convection --- the dominant regime for slowly circulated water-cooled PbLi
concepts --- is not yet modelled; the present tool addresses self-cooled and
dual-coolant concepts where forced-flow pressure drop dominates.

The pipeline is nevertheless sufficient to answer the design question it targets:
whether a proposed channel topology on a given stellarator equilibrium is feasible,
and at what pumping cost the required TBR is reached.

\section{Conclusion}

\texttt{mhd-tbr-stell} couples, for the first time in public code, the MHD
pressure drop and the tritium breeding ratio of liquid-metal flow channels
conforming to stellarator magnetic surfaces. The tool chain --- VMEC Fourier
geometry, correlation-based duct network, epotFoam anchoring, OpenMC neutronics ---
is assembled entirely from validated open-source components, and every mock result
in this draft is scheduled to be replaced by production computations on a W7-X--like
and an LHD-like configuration. Source code is available at
\url{https://github.com/lzpel/mhd-tbr-stell}.

\bibliographystyle{plain}
\bibliography{refs}

\end{document}
